\documentclass[11pt]{article}
\usepackage[margin=1.05in]{geometry}
\usepackage{amsmath, amssymb}
\usepackage{booktabs}
\usepackage{multirow}
\usepackage{xcolor}
\usepackage{hyperref}
\hypersetup{colorlinks=true, linkcolor=blue!50!black, urlcolor=blue!50!black}

\newcommand{\code}[1]{\texttt{#1}}
\DeclareMathOperator{\hn}{HN}

% Table bodies are generated by
% code/analysis/time_block_forecast/simstudy/expA_breakdown.R and are
% never edited by hand; this report transcribes no number.
\newcommand{\tblpath}{../../code/analysis/time_block_forecast/simstudy/output/tables}
\newcommand{\tbl}[1]{\input{\tblpath/#1_body.tex}}

\title{\textbf{Per-block threshold Brier scores for the time-block
forecast}\\[2pt]
\large Per-block and per-dataset breakdowns, and why the rule was
demoted}
\author{Score report}
\date{4 August 2026}

\begin{document}
\maketitle

\begin{abstract}
This is the companion to \code{tex/timeblock\_brier\_fullseries}. It
scores the \emph{same fits}, on the same ten datasets, five blocks and
three methods, and differs in exactly one respect: the Brier thresholds
are the validation block's own quantiles rather than the dataset's full
series. Because the two reports share a cache, the comparison between
them is controlled --- no difference between them can be attributed to
different fits, samples, or priors. That comparison is the point of this
report, and it is stronger than the look-ahead objection that originally
motivated the demotion. The per-block rule does not shrink the AR(2)
advantage; it \emph{inflates the noise around it}. At setting~5,
$q_{95}$, its estimated gap is $-18.1 \times 10^{-3}$, more than twice
the full-series estimate of $-8.9 \times 10^{-3}$, and yet it reaches
$p = 0.110$ where the full-series rule reaches $p = 0.002$. Per-dataset
agreement falls from $10/10$ to $7/10$. The mechanism is measured
directly in \S\ref{sec:pinning}: the rule holds the realised exceedance
rate at exactly $1-p$ in all fifty (dataset, block) cells --- standard
deviation zero --- so the threshold is a random statistic of the very
window being scored, and the per-block randomness it injects survives the
pairing. A further symptom is that this rule's estimated gap is nearly
\emph{constant} in $p$, where the full-series gap decays with the base
rate as it should.
\end{abstract}

\section{Scope and design}

Identical to the sibling report: the Experiment~A campaign from
\code{simstudy/output/results/scores\{5,7\}\_hn.RData}, seams at
$T_o = 50, 80, 110, 140, 170$, $H = 15$, three methods (i.i.d., AR(1),
AR(2)), ten datasets, five blocks, $\lambda \sim \hn(0, 20)$,
\code{hn\_prior = TRUE}, no missing cells and hence $n = 10$ with no
clean-$n$ selection. Settings 5 (near-unit-root AR(2),
$\phi = (0.15, 0.80)$) and 7 (ARFIMA$(0, 0.45, 0)$, Hurst $0.95$).

The one difference: every table below reads \code{brier.lead.blockq}
instead of \code{brier.lead}. \code{scores.R} computes both arrays in the
same pass, from the same predictive draws, so the two reports are two
scorings of one campaign rather than two campaigns.

Thresholds $q_{90}$ and $q_{95}$ are co-primary; the eleven-point grid is
in the appendix. Brier only --- CRPS is invariant to the threshold rule
and so has nothing to say about the distinction this report exists to
draw. The retired campaign's CRPS results are in
\code{tex/timeblock\_expABC\_legacy}.

\section{The per-block threshold rule}

For each block the thresholds are taken from the held-out window itself,
\[
  \tau_p^{\mathrm{blk}}(d,s,b)
  \;=\; \mathrm{quantile}\bigl(y_{\mathrm{val}}(d,s,b),\,p\bigr),
\]
recomputed for every (dataset, setting, block) and shared across methods
and leads within that block. It was the rule in force for every Brier
table written before 2026-07-31, and it is retained in the caches as
\code{brier.lead.blockq}, tagged so that
\code{brier\_threshold\_basis} distinguishes it and the merge gate
refuses to average the two definitions together.

Two objections were raised against it. The first is that it is
look-ahead: the threshold is a statistic of the window being forecast, so
the scoring rule is not implementable in a real forecast setting. That
objection is sound but qualitative. The second is that it pins the
exceedance base rate, so a level error cannot move it.
\S\ref{sec:pinning} measures the pinning, and \S\ref{sec:contrast}
measures what it costs.

\section{Per-block breakdown}

\tbl{expA_block_blockq}

\noindent The levels are the tell. Under this rule setting~5's i.i.d.\
scores at $q_{95}$ run from $0.0551$ to $0.1221$ across the five blocks
--- a factor of $2.2$ between the easiest and hardest window. Under the
full-series rule the same five blocks run from $0.0154$ to $0.1060$, a
factor of $6.9$: three times the dynamic range, on identical forecasts.
The rule compresses the
blocks toward each other because it gives each block a threshold matched
to its own turbulence: a quiet window is no longer scored as quiet. The
block that most disagrees between the two rules is block~3, the quietest
one, whose i.i.d.\ level rises from $0.0154$ to $0.0789$.

AR(2) still leads i.i.d.\ in all five blocks at setting~5, and in two of
five at setting~7 --- one fewer than the full-series rule reports.

\section{Per-dataset breakdown}

\tbl{expA_dataset_blockq}

\noindent At setting~5 AR(2) wins $7/10$ datasets at both thresholds,
against $10/10$ ($q_{90}$) and $9/10$ ($q_{95}$) under the full-series
rule --- three datasets change side purely from the threshold
definition, on identical forecasts. This is the per-dataset face of the
variance inflation quantified next.

\section{Paired tests}

\tbl{expA_tests_blockq}

\noindent Same construction as the sibling report: one-sided, dataset
unit, $n = 10$, lead window averaged before blocks. Only one contrast
attains significance anywhere --- AR(2)$-$AR(1) over leads $1$--$15$,
at both thresholds, marginally ($t\,p = 0.046$ and $0.039$). The
AR(2)$-$i.i.d.\ contrast, which the full-series rule resolves at
$p \le 0.004$ everywhere, does not separate at all here. Setting~7 is
null throughout, as it is under the other rule.

\section{The two rules on identical fits}
\label{sec:contrast}

\tbl{expA_rulecontrast}

\noindent This is the controlled comparison the retired campaign could
never make, because its Brier tables and its successors came from
different fits. Here the fits, the datasets, the blocks and the
predictive draws are the same objects in memory; only the threshold
differs.

\paragraph{The rule does not shrink the signal.} At setting~5 the
per-block rule's point estimates are \emph{larger} in magnitude than the
full-series rule's at $q_{95}$ --- $-18.1$ against $-8.9$ (in $10^{-3}$)
for AR(2)$-$i.i.d., and $-21.5$ against $-10.5$ for AR(2)$-$AR(1). If
the demotion cost only sensitivity, this is not what one would see.

\paragraph{It inflates the noise.} Despite the larger gaps, the
per-block rule cannot separate: $p = 0.110$ against $0.002$ for
AR(2)$-$i.i.d.\ at $q_{95}$. A larger estimate with a much larger
$p$-value means the across-dataset variance has grown faster than the
mean. The pairing is what suffers: two methods scored in the same block
still share a threshold, but that threshold is itself a draw, and its
randomness enters every dataset's difference differently.

\paragraph{It is nearly blind to the threshold.} The appendix sweeps
both rules across the eleven-point grid. Under the full-series rule
setting~5's gap decays smoothly with $p$, from $-25.8$ at $q_{90}$ to
$-0.6$ at $q_{995}$, tracking the base rate. Under the per-block rule it
is almost flat --- $-19.5$ at $q_{90}$, $-11.7$ at $q_{995}$ --- and its
$p$-value barely moves off $0.11$ across the entire grid. A scoring rule
whose answer does not depend on which tail you ask about is not
measuring the tail; it is measuring the fixed $1-p$ of events it has
guaranteed itself.

\section{Base-rate pinning, measured}
\label{sec:pinning}

The claim that this rule pins the exceedance rate has been asserted
before but not shown. It is a statement about the data alone --- no fit
enters --- so it can be checked directly. For each of the fifty
(dataset, block) cells the realised exceedance rate of the validation
window is computed under both rules.

\tbl{expA_pinning}

\noindent The pinning is exact. Under the per-block rule every one of the
fifty cells realises a rate of precisely $1-p$: mean $0.1000$ at
$q_{90}$ and $0.0500$ at $q_{95}$, with standard deviation \emph{zero}
and identical minimum and maximum. Under the full-series rule the same
cells range from $0$ to $0.9306$ at setting~5, $q_{90}$.

This is the mechanism behind both preceding sections. A window that ran
far above its training level cannot register as unusual, because its
threshold rose with it; a window in which nothing happened is handed
exactly as many ``exceedances'' as a violent one. The base rate carries
no information, so the only channel left to the score is the shape of the
predictive around a moving target --- which is both a weaker signal and a
noisier one.

\section{Provenance}

Caches, scripts and gates as in the sibling report:
\code{simstudy/output/results/scores\{5,7\}\_hn.RData},
\code{brier\_threshold\_basis = "full\_series"} recording the primary
rule, \code{hn\_prior = TRUE}; tables from
\code{Rscript expA\_breakdown.R --prior=hn --settings=5,7}, which emits
both rules' fragments from one load and recomputes the pinning table from
\code{simdata.RData}. No number in this document is transcribed by hand.

The \code{\_hn} suffix is the lambda-prior arm, which every per-arm
artifact has carried since 2026-08-05. Both this report and its sibling
are the HN arm; the threshold rule is the only thing that differs between
them.

Sibling: \code{tex/timeblock\_brier\_fullseries}. Frozen record of the
retired campaign, settings 4 and 6, Experiments~B and~C, and all CRPS
results: \code{tex/timeblock\_expABC\_legacy}.

\appendix
\section{The full threshold grid}

\tbl{expA_qsweep_blockq}

\noindent Per-block rule, pooled over leads $1$--$15$ at the dataset
unit. Compare the full-series sweep in the sibling report's appendix: the
gap here is nearly constant in $p$ and the $p$-value never leaves the
neighbourhood of $0.11$, while the dataset and block counts are frozen at
$7/10$ and $5/5$ across all eleven thresholds --- the signature of a
score whose base rate has been fixed by construction.

\end{document}
